\documentclass[10pt,letterpaper]{article}
\usepackage[T1]{fontenc}
\usepackage[utf8]{inputenc}
\usepackage{lmodern}
\usepackage[letterpaper,margin=0.84in,headheight=15pt]{geometry}
\usepackage[table]{xcolor}
\definecolor{DeepNavy}{HTML}{173450}
\definecolor{Teal}{HTML}{087F8C}
\definecolor{Slate}{HTML}{536578}
\definecolor{SoftBlue}{HTML}{F1F6FA}
\definecolor{SoftGray}{HTML}{F6F8FA}
\definecolor{Crimson}{HTML}{B34843}
\usepackage{microtype}
\setlength{\emergencystretch}{1.5em}
\usepackage{booktabs,longtable,tabularx,array}
\usepackage{amssymb}
\usepackage{ragged2e}
\usepackage{enumitem}
\usepackage{fancyhdr,lastpage}
\usepackage{titlesec}
\usepackage{listings}
\usepackage[most]{tcolorbox}
\usepackage{needspace}
\usepackage{hyperref}
\hypersetup{colorlinks=true,linkcolor=DeepNavy,urlcolor=Teal,citecolor=DeepNavy,pdfauthor={Software Engineering Technical Documentation}}
\setlength{\parindent}{0pt}
\setlength{\parskip}{5pt plus 1pt}
\setlist[itemize]{itemsep=2pt,topsep=3pt,leftmargin=1.55em}
\setlist[enumerate]{itemsep=2pt,topsep=3pt,leftmargin=1.7em}
\titleformat{\section}{\Large\bfseries\color{DeepNavy}}{\thesection}{0.75em}{}[\vspace{-3pt}\color{Teal}\titlerule]
\titleformat{\subsection}{\large\bfseries\color{DeepNavy}}{\thesubsection}{0.65em}{}
\titleformat{\subsubsection}{\normalsize\bfseries\color{Teal}}{\thesubsubsection}{0.55em}{}
\pagestyle{fancy}
\fancyhf{}
\fancyhead[L]{\sffamily\footnotesize MISRA C:2023 | Engineering Documentation}
\fancyhead[R]{\sffamily\footnotesize\color{Slate} Independent companion}
\fancyfoot[L]{\sffamily\scriptsize Original explanatory material; consult licensed MISRA standard}
\fancyfoot[R]{\sffamily\scriptsize\thepage\ / \pageref{LastPage}}
\renewcommand{\headrulewidth}{0.3pt}
\newcommand{\ID}[1]{\textsf{\bfseries #1}}
\newcommand{\SourceNote}[1]{\par\smallskip{\footnotesize\color{Slate}\textit{Trace: #1.}}\par}
\newtcolorbox{keybox}[1][]{enhanced,breakable,colback=SoftBlue,colframe=Teal,boxrule=0.5pt,arc=1mm,title={Key point},fonttitle=\bfseries,#1}
\newtcolorbox{riskbox}[1][]{enhanced,breakable,colback=SoftGray,colframe=Crimson,boxrule=0.5pt,arc=1mm,title={Verification boundary},fonttitle=\bfseries,#1}
\lstdefinestyle{engineerc}{language=C,basicstyle=\ttfamily\footnotesize,keywordstyle=\color{DeepNavy}\bfseries,commentstyle=\color{Slate}\itshape,stringstyle=\color{Teal},numbers=left,numberstyle=\tiny\color{Slate},stepnumber=1,numbersep=8pt,breaklines=true,breakatwhitespace=false,frame=single,framesep=5pt,rulecolor=\color{Slate!25},backgroundcolor=\color{SoftGray},tabsize=2,showstringspaces=false,columns=fullflexible,keepspaces=true}
\lstdefinestyle{plaincode}{basicstyle=\ttfamily\footnotesize,numbers=none,breaklines=true,frame=single,rulecolor=\color{Slate!25},backgroundcolor=\color{SoftGray},columns=fullflexible,keepspaces=true}
\newcommand{\TrainingTag}[1]{\textsf{\footnotesize\color{Teal}\bfseries #1}}
\newcommand{\DocNotes}{\begin{keybox}This publication is an independently written educational and operational companion. It is \emph{not} the licensed MISRA guideline text, a MISRA-endorsed tool, or a certification of project compliance. Identifiers and published classifications are used for traceability; check exact requirement wording, amplification, applicability, and exceptions in the licensed standard.\end{keybox}}
\AtBeginDocument{\color{black}}

\title{\vspace{-2.3em}\color{DeepNavy}\Huge\bfseries Deviation Management\\[0.45em]\large\normalfont\color{Teal} Formal risk decisions, governance and controlled examples}
\author{Independent Software Engineering Reference Series}
\date{Edition 1.0 \quad | \quad October 2026}
\begin{document}
\maketitle
\vspace{-0.4em}
\begin{center}\color{Slate}\rule{0.83\linewidth}{0.4pt}\end{center}
\DocNotes
\tableofcontents
\newpage

\section{Disposition authority}
For MISRA C:2023, Mandatory guidelines cannot be formally deviated from; Required guidelines may be deviated from through the compliance process; Advisory departures require appropriate documentation even when a formal deviation is not necessary. Approved local re-categorization can strengthen enforcement. The exact conditions of a compliant project claim are governed by MISRA Compliance:2020.
\section{Decision process}
\begin{enumerate}
\item \textbf{Classify}: identify exact guideline, edition, category, exception and all affected configurations.
\item \textbf{Confirm}: replicate the diagnostic, assess cross-file context and verify it is an actual violation.
\item \textbf{Correct by default}: prefer code or design repair unless the repair creates disproportionate technical risk or violates an explicit hardware contract.
\item \textbf{Check eligibility}: reject deviations for Mandatory items; escalate unsupported exceptions to scope/design change.
\item \textbf{Evaluate}: analyze failure mechanism, use-case boundaries, plausible consequences, likelihood, detectability and residual risk.
\item \textbf{Mitigate}: add contract assertions, test cases, static proofs, isolation boundaries, restrictions or external review.
\item \textbf{Approve}: secure authorized independent safety/security and technical signoff before merge/release.
\item \textbf{Control}: register affected locations, revision, owner, expiry trigger, permitted use-case and evidence.
\item \textbf{Revalidate}: expire or review on compiler change, target change, refactor, incident, new analyzer result, or release milestone.
\end{enumerate}
\section{Deviation record field catalog}
\begin{tabularx}{\linewidth}{@{}p{.26\linewidth}X@{}}\toprule
Field & Minimum engineering rationale\\\midrule
Record ID, status, owner & Immutable identifier, draft/approved/rejected/expired and accountable developer\\
Rule ID and category & Exact applicable licensed edition and category; no Mandatory waivers\\
Scope and locations & Repository hash, files, relevant functions, configurations and affected products\\
Purpose & Engineering problem requiring deviation and why alternatives are unsuitable\\
Mechanism and impact & Behavior at issue, safety/security hazard, fault propagation and downstream exposure\\
Compensating controls & Encapsulation, tests, proof obligations, hardware assumptions and monitoring\\
Evidence & Reproducible findings, compiler info, test results, reviewer notes and related requirements\\
Review and expiry & Approval identities, date, release validity and change-triggered reassessment\\\bottomrule\end{tabularx}
\section{Example: memory-mapped peripheral address}
\begin{keybox}\textbf{Illustrative only}: A project may need a documented exception for a tightly controlled integer-to-object-pointer conversion (check the exact treatment of Rule 11.4 and related rules in the licensed publication). The demonstration below is not a blanket approval or proof of safety.\end{keybox}
\begin{lstlisting}[style=engineerc]
/* Project-specific device address: hypothetical target only. */
#define DEVICE_BASE_ADDR UINT32_C(0x40002000)
static volatile uint32_t * const device_status =
    (volatile uint32_t *)(uintptr_t)DEVICE_BASE_ADDR;
\end{lstlisting}
\textbf{Risk rationale:} The C implementation and target memory map must support the mapping, alignment, width, volatility and register semantics. Encapsulate the conversion in one reviewed module. Reject use on a host build or a different target. Require a trace to the silicon manual, linker map and access tests.
\textbf{Approval example:} ID DEV-EX-0001, Rule 11.4, target rev A only, single source location, reviewed by hardware owner plus assurance lead, expiry on target or compiler revision, with explicit prohibition on casual reuse.
\section{Deviation permits and records}
A \emph{deviation permit} describes an approved recurring use-case and review requirements; a \emph{deviation record} applies it to identifiable project violations. A permit does not automatically approve any particular violation. Maintain location-level applicability and a controlled register. See Section 4 of MISRA Compliance:2020 for the authority model and required content.
\section{Anti-patterns}
\begin{itemize}
\item ``False positive'' selected without explaining checker limitations or actual semantics.
\item Permanent project-wide suppressions with no affected-location inventory.
\item Treating tool comments as risk approval.
\item Waiving a Mandatory guideline.
\item Approving a deviation after production release or ignoring cross-configuration variants.
\item Allowing a baseline count to grow without owner, budget, or scheduled closure.
\end{itemize}
\section{Deviation acceptance checklist}
\begin{itemize}
\item[$\square$] Actual guideline text and exceptions consulted.
\item[$\square$] Rule noncompliance independently verified.
\item[$\square$] Deviation eligibility proven and project category checked.
\item[$\square$] Alternatives and risks assessed and compensating controls demonstrated.
\item[$\square$] Locations, target, configuration and release scope captured.
\item[$\square$] Approvals, revision, expiry and re-review trigger recorded.
\end{itemize}

\end{document}
